\documentclass[12pt,a4paper,psamsfonts]{amsart}

% Layout and theorem styles follow arXiv:2309.07005v1.
\usepackage{amssymb,amscd,amsxtra,calc}
\usepackage{mathrsfs,mathtools,cmmib57}
\usepackage{mathabx}
\usepackage[colorlinks,linkcolor=blue,anchorcolor=blue,citecolor=green]{hyperref}

\setlength{\topmargin}{0cm}
\setlength{\oddsidemargin}{0cm}
\setlength{\evensidemargin}{0cm}
\setlength{\marginparwidth}{0cm}
\setlength{\marginparsep}{0cm}
\setlength{\textheight}{\paperheight - 2in -35pt}
\setlength{\textwidth}{\paperwidth - 2in}
\setlength{\headheight}{12.5pt}
\setlength{\headsep}{25pt}
\setlength{\footskip}{30pt}
\renewcommand{\baselinestretch}{1.3}
\pagestyle{headings}
\setlength{\emergencystretch}{2em}

\theoremstyle{plain}
\newtheorem{theorem}{Theorem}[section]
\newtheorem{lemma}[theorem]{Lemma}
\theoremstyle{definition}
\newtheorem{remark}[theorem]{Remark}

\newcommand{\R}{\mathbb R}
\newcommand{\Q}{\mathbb Q}
\newcommand{\OO}{\mathcal O}
\DeclareMathOperator{\Supp}{Supp}
\DeclareMathOperator{\codim}{codim}
\DeclareMathOperator{\coeff}{coeff}

\hypersetup{pdftitle={Negativity Lemma}}

\title[Negativity Lemma]{Negativity Lemma}
\author{}
\date{}

\begin{document}
\maketitle
\markboth{NEGATIVITY LEMMA}{NEGATIVITY LEMMA}

\section{Detailed argument}

\begin{lemma}\label{lem:antiample}
Let $f\colon X\to Y$ be a projective birational morphism between
normal varieties, with $Y$ affine. Then there exists an effective
Cartier divisor $E$ on $X$ such that $-E$ is $f$-ample.
\end{lemma}

\begin{proof}
Since $f$ is projective, we can choose an $f$-ample Cartier divisor
$D$ on $X$. Put $\mathcal F=f_*\OO_X(-D)$. Since $f$ is projective,
$\mathcal F$ is coherent. Since $f$ is birational, it is an
isomorphism over a dense open subset of $Y$, on which $\mathcal F$
is invertible. Thus $\mathcal F$ has generic rank one; in particular,
$\mathcal F\neq0$.

Since $Y$ is affine, a nonzero quasi-coherent sheaf on $Y$ has a
nonzero global section. Hence we may choose
\[
0\neq s\in H^0(Y,\mathcal F)=H^0(X,\OO_X(-D)).
\]
Viewing $\OO_X(-D)$ as a subsheaf of the sheaf of rational functions,
we regard $s$ as a nonzero rational function on $X$. Set
\[
E=-D+\operatorname{div}(s).
\]
The condition $s\in H^0(X,\OO_X(-D))$ says precisely that $E$ is
an effective Cartier divisor. Indeed, if $d$ is a local equation
of $D$, then $s/d$ is a nonzero regular function and is a local
equation of $E$. Finally, $-E\sim D$, so $-E$ is $f$-ample.
\end{proof}

\begin{lemma}\label{lem:exceptional-support}
Let $f\colon X\to Y$ be a projective birational morphism between
varieties, with $Y$ normal, and let $E$ be an effective Cartier
divisor on $X$ such that $-E$ is $f$-ample. Then $\Supp E$ contains
the exceptional locus $\operatorname{Exc}(f)$.
\end{lemma}

\begin{proof}
Since $Y$ is normal, $\operatorname{Exc}(f)$ is covered by complete
integral curves contracted by $f$. For every such curve $C$, the
$f$-ampleness of $-E$ gives $E\cdot C<0$. Since $E$ is effective,
this forces $C\subseteq\Supp E$. Hence
$\operatorname{Exc}(f)\subseteq\Supp E$.
\end{proof}

\begin{theorem}[Projective negativity lemma]\label{thm:negativity}
Let $f\colon X\to Y$ be a projective birational morphism between normal
varieties. Let $D$ be an $\R$-Cartier divisor on $X$ such that $-D$ is
nef over $Y$. Then:
\begin{enumerate}
\item $D$ is effective if and only if $f_*D$ is effective.
\item If $D$ is effective, then for every $y\in Y$, either
\[
f^{-1}(y)\cap\Supp D=\varnothing
\qquad\text{or}\qquad
f^{-1}(y)\subset\Supp D.
\]
\end{enumerate}
\end{theorem}

\begin{proof}[Proof of Theorem~\ref{thm:negativity}]
Both assertions are local on $Y$, so we may assume that $Y$ is affine.

\textup{(1)}
One direction is trivial. For the other direction, assume $f_*D\geq0$.
Since $f$ is birational, every prime divisor with negative coefficient
in $D$ is $f$-exceptional.

Suppose, for a contradiction, that $D$ is not effective. Choose an effective
Cartier divisor $E$ on $X$ such that $-E$ is $f$-ample, as in
Lemma~\ref{lem:antiample}. By Lemma~\ref{lem:exceptional-support},
the coefficient $\coeff_P(E)>0$ whenever $\coeff_P(D)<0$. Set
\[
e:=\max_{\coeff_P(D)<0}\frac{-\coeff_P(D)}{\coeff_P(E)}>0.
\]
Then $D+eE\geq0$, and we can find an $f$-exceptional prime
divisor $F$ such that $F\not\subset\Supp(D+eE)$.
Choose a closed point $x\in F\setminus\Supp(D+eE)$ and a
curve $C\subset F$ through $x$ contracted by $f$.
Then $C\not\subset\Supp(D+eE)$, so $(D+eE)\cdot C\geq0$.
On the other hand,
\[
(D+eE)\cdot C=D\cdot C+e(E\cdot C)<0,
\]
because $-D$ is $f$-nef, $-E$ is $f$-ample, and $e>0$. This is a
contradiction, so $D\geq0$.

\textup{(2)}
Assume $D\geq0$. Since $f$ is projective and birational and $Y$ is
normal, its fibers are connected by Stein factorization.
Suppose $f^{-1}(y)\cap\Supp D\neq\varnothing$. If
$f^{-1}(y)\not\subset\Supp D$, connectedness gives a curve
$C\subset f^{-1}(y)$ meeting $\Supp D$ but not contained in it.
Then $D\cdot C>0$, contradicting the $f$-nefness of $-D$.
Thus $f^{-1}(y)\subset\Supp D$.
\end{proof}

\begin{theorem}[Proper negativity lemma]\label{thm:proper-negativity}
Let $f\colon X\to Y$ be a proper birational morphism between normal
varieties. Let $D$ be an $\R$-Cartier divisor on $X$ such that $-D$ is
nef over $Y$. Then:
\begin{enumerate}
\item $D$ is effective if and only if $f_*D$ is effective.
\item If $D$ is effective, then for every $y\in Y$, either
\[
f^{-1}(y)\cap\Supp D=\varnothing
\qquad\text{or}\qquad
f^{-1}(y)\subset\Supp D.
\]
\end{enumerate}
\end{theorem}

\begin{proof}
By Chow's lemma and normalization, choose a proper birational
morphism $\pi\colon X'\to X$ with $X'$ normal such that
$f'=f\circ\pi\colon X'\to Y$ is projective. Set $D'=\pi^*D$.
By the projection formula, $-D'$ is $f'$-nef, and
\[
\pi_*D'=D,\qquad f'_*D'=f_*D.
\]
If $f_*D\geq0$, Theorem~\ref{thm:negativity}\textup{(1)} gives
$D'\geq0$, hence $D=\pi_*D'\geq0$. The converse is immediate.

If $D\geq0$, then $D'\geq0$ and
$\Supp D'=\pi^{-1}(\Supp D)$. By
Theorem~\ref{thm:negativity}\textup{(2)}, each fiber of $f'$ is
either disjoint from $\Supp D'$ or contained in it. Since
$f'^{-1}(y)\to f^{-1}(y)$ is surjective, the same holds for $f$
and $D$.
\end{proof}

\section{A concise proof}

\begin{proof}[Proof of Theorem~\ref{thm:proper-negativity}]
We may work over an affine open subset of $Y$. By Chow's lemma
and normalization, choose $\pi\colon X'\to X$ proper birational
with $X'$ normal and $f\circ\pi$ projective. The projection formula
preserves relative nefness, and
\[
\pi_*\pi^*D=D,\qquad
\Supp(\pi^*D)=\pi^{-1}(\Supp D)\quad\text{if }D\geq0.
\]
Thus both assertions descend from $X'$ to $X$, and we may assume
that $f$ is projective.

\textup{(1)} Suppose $f_*D\geq0$ but $D\not\geq0$.
Choose an $f$-ample Cartier divisor $A$. The coherent sheaf
$f_*\OO_X(-A)$ has generic rank one and hence, since $Y$ is affine,
a nonzero global section $s$. Then $E=-A+\operatorname{div}(s)\geq0$
and $-E$ is $f$-ample. Every contracted curve $C$ satisfies
$E\cdot C<0$, so $C\subset\Supp E$. Such curves cover
$\operatorname{Exc}(f)$; hence $\operatorname{Exc}(f)\subset\Supp E$.

All negative components of $D$ are $f$-exceptional. Set
\[
e:=\max_{\coeff_P(D)<0}
\frac{-\coeff_P(D)}{\coeff_P(E)}>0.
\]
Then $D+eE\geq0$, and some $f$-exceptional prime divisor $F$ is
not contained in $\Supp(D+eE)$. Choose a contracted curve
$C\subset F$ not contained in this support. We obtain
\[
0\leq(D+eE)\cdot C=D\cdot C+e(E\cdot C)<0,
\]
a contradiction. The reverse implication is immediate.

\textup{(2)} Assume $D\geq0$ and
$f^{-1}(y)\cap\Supp D\neq\varnothing$. The fiber is connected by
Stein factorization. If it is not contained in $\Supp D$, there is
a curve in the fiber meeting $\Supp D$ but not contained in it.
Its intersection with $D$ is positive, contrary to the
$f$-nefness of $-D$. Thus $f^{-1}(y)\subset\Supp D$.
\end{proof}

\section{Extension to quasi-projective bases}

Lemma~\ref{lem:antiample} remains valid when the affine hypothesis
on $Y$ is replaced by quasi-pro\-jec\-tivity: for a projective birational
morphism $f\colon X\to Y$ between normal varieties with $Y$
quasi-projective, there is an effective Cartier divisor $E$ such
that $-E$ is $f$-ample. This strengthening also requires no
Riemann--Roch.

\begin{proof}
Choose an $f$-ample line bundle $\mathcal L$ and put
$\mathcal F=f_*\mathcal L^{-1}$. As in
Lemma~\ref{lem:antiample}, $\mathcal F$ is coherent of generic rank
one, so $\mathcal F\neq0$.

Choose a very ample Cartier divisor $H$
on $Y$. By the global-generation criterion for ampleness,
$\mathcal F\otimes\OO_Y(mH)$ is generated by global sections for
$m\gg0$. It is nonzero, so it has a nonzero global section.
The projection formula gives
\[
\begin{split}
0\neq s&\in H^0\bigl(Y,\mathcal F\otimes\OO_Y(mH)\bigr)\\
&=H^0\bigl(X,\mathcal L^{-1}\otimes\OO_X(mf^*H)\bigr).
\end{split}
\]
Its zero divisor $E$ satisfies
\[
\OO_X(-E)\simeq\mathcal L\otimes f^*\OO_Y(-mH).
\]
Tensoring with a line bundle pulled back from the base preserves
relative ampleness, so $-E$ is again $f$-ample. If $Y$ is projective,
then $X$ is projective. Taking $\mathcal L=\OO_X(A)$ for
any very ample divisor $A$ on $X$ gives $E\sim mf^*H-A$ and hence
the decomposition
\[
mf^*H\sim A+E,\qquad -E\sim_f A.
\]


Lemma~\ref{lem:exceptional-support} further shows
that every exceptional prime divisor lies in $\Supp E$.
The only additional input is global generation after a sufficiently
large ample twist. No Riemann--Roch or section-growth estimate is
needed.
\end{proof}

\begin{remark}
The affine hypothesis is essential to the \emph{untwisted argument},
not to the conclusion over a quasi-projective base. A nonzero
coherent sheaf on a projective variety need not have any global
sections. For example, with $f=\operatorname{id}_{\mathbf P^1}$ and
$\mathcal L=\OO_{\mathbf P^1}(1)$, one has
$f_*\mathcal L^{-1}=\OO_{\mathbf P^1}(-1)\neq0$ but
$H^0(X,\mathcal L^{-1})=0$. The twist by $mH$ repairs precisely
this failure. On an arbitrary normal base, a global very ample
$H$ need not exist; for the negativity lemma it suffices to apply
the affine construction on each affine open subset of the base.
\end{remark}

\begin{remark}
For comparison, the projective decomposition also has the following proof, making
explicit the section-growth argument usually used in Kodaira's
lemma. Assume that $Y$ is projective, put $n=\dim X=\dim Y\geq1$,
choose very ample Cartier divisors $H$ on $Y$ and $A$ on $X$,
and set $M=f^*H$. Since $Y$ is normal and $f$ is proper birational,
$f_*\OO_X=\OO_Y$. Thus
\[
h^0(X,\OO_X(mM))=h^0(Y,\OO_Y(mH))
=\frac{H^n}{n!}m^n+O(m^{n-1}),\qquad H^n>0.
\]
The last equality follows from Serre vanishing and the Hilbert
polynomial of the very ample divisor $H$. In particular, $M$ is
big; it need not be ample.

Choose an effective member $T\in|A|$. It is a Cartier divisor on
the integral variety $X$, and $\dim T=n-1$. The coherent sheaf
$\mathcal G=f_*\OO_T$ is supported on $f(T)$, whose dimension is at
most $n-1$. By the projection formula and the Hilbert polynomial
of $\mathcal G$ on $Y$,
\[
\begin{split}
h^0(T,\OO_T(mM))
&=h^0\bigl(Y,\mathcal G\otimes\OO_Y(mH)\bigr)\\
&=O(m^{n-1}).
\end{split}
\]
No normality or irreducibility of $T$ is needed. Taking global
sections of the exact sequence
\[
0\longrightarrow\OO_X(mM-T)\longrightarrow\OO_X(mM)
\longrightarrow\OO_T(mM)\longrightarrow0
\]
gives
\[
\begin{split}
h^0(X,\OO_X(mM-A))
&\geq h^0(X,\OO_X(mM))-h^0(T,\OO_T(mM))\\
&\geq \frac{H^n}{n!}m^n-O(m^{n-1})>0
\end{split}
\]
for $m\gg0$. We only use left exactness here; surjectivity of the
restriction map is unnecessary. A nonzero section therefore gives
an effective Cartier divisor $E\sim mM-A$. Equivalently,
\[
f^*H\sim_{\Q}\frac1m A+\frac1m E,
\qquad -E\sim_f A.
\]
This is the requested decomposition into an ample divisor and an
effective divisor. Lemma~\ref{lem:exceptional-support} shows that
$\Supp E$ contains every exceptional prime divisor. When $n=0$,
the assertion is immediate with $E=0$.

The essential numerical input in this second proof is that the
space of sections on $X$ grows like $m^n$, whereas the space of
restrictions to $T$ grows at most like $m^{n-1}$. It is often
expressed as an asymptotic Riemann--Roch argument. However, full
Riemann--Roch is \emph{not logically necessary}: the Hilbert
polynomial and Serre vanishing on $Y$ already give both estimates,
and the first proof needs only global generation after an ample
twist. Riemann--Roch alone computes an Euler characteristic, not
$h^0$; a proof invoking it must also justify the relevant
cohomological estimates. In particular, one cannot apply Serre
vanishing on $X$ merely because $f^*H$ is big and nef. Neither
construction uses the Hodge index theorem, and $E$ may have
nonexceptional components.
\end{remark}

\end{document}
